\documentclass[conference]{IEEEtran}
\IEEEoverridecommandlockouts
\usepackage{cite}
\usepackage{amsmath,amssymb,amsfonts}
\usepackage{algorithmic}
\usepackage{algorithm}
\usepackage{graphicx}
\usepackage{textcomp}
\usepackage{xcolor}
\usepackage{url}
\def\BibTeX{{\rm B\kern-.05em{\sc i\kern-.025em b}\kern-.08em
    T\kern-.1667em\lower.7ex\hbox{E}\kern-.125emX}}
\begin{document}

\title{AI-Powered Agentic Hiring Platform: A Multi-Agent System for Intelligent Recruitment, Explainable Job Matching, and Candidate Development}

\author{
\IEEEauthorblockN{Mrs.~Joby Anu Mathew}
\IEEEauthorblockA{\textit{Assoc. Prof, Dept.\ of AI \& Data Science} \\
\textit{Viswajyothi College of Engineering and}\\
Technology Vazhakulam, India \\
jobyanu@vjcet.org}

\and
\IEEEauthorblockN{Abhijith Binosh}
\IEEEauthorblockA{\textit{Dept.\ of AI \& Data Science} \\
\textit{Viswajyothi College of Engineering and }\\
Technology Vazhakulam, India \\
abhijith22ra124@vjcet.org}

\and
\IEEEauthorblockN{Jose Martin Binu}
\IEEEauthorblockA{\textit{Dept.\ of AI \& Data Science} \\
\textit{Viswajyothi College of Engineering and}\\
Technology Vazhakulam, India \\
jose22ra239@vjcet.org}

\and

\hspace{0.6cm}
\IEEEauthorblockN{R.~Jayadev}
\IEEEauthorblockA{\textit{Dept.\ of AI \& Data Science} \\
\textit{Viswajyothi College of Engineering and}\\
 Technology Vazhakulam, India \\
jayadev22ra125@vjcet.org}

\and
\IEEEauthorblockN{Suryanarayanan N.~J.}
\IEEEauthorblockA{\textit{Dept.\ of AI \& Data Science} \\
\textit{Viswajyothi College of Engineering and}\\
 Technology Vazhakulam, India \\
surya22ra129@vjcet.org}
}

\maketitle

% ─────────────────────────────────────────────────────────────
\begin{abstract}
The recruitment industry faces significant challenges in efficiently matching candidates with job opportunities due to high application volumes, subjective evaluation criteria, and fragmented tooling. This paper presents an AI-Powered Agentic Hiring Platform that integrates a multi-agent architecture, hybrid AI/NLP-based candidate shortlisting, automated interview scheduling, and a suite of candidate development tools. The system's novel contributions include: (1)~a hybrid shortlisting agent combining TF-IDF cosine similarity with Google Gemini~AI for explainable match scoring; (2)~a break-aware round-robin scheduling algorithm with Fisher-Yates randomisation guaranteeing $\Delta_{\text{load}}\!\leq\!1$; (3)~a remote coding assessment module backed by the Piston API; (4)~a Gemini-powered resume optimiser with hallucination guard and PDF export; (5)~a career roadmap generator using Groq Llama~3.3-70B; (6)~a profile-driven job recommender aggregating from Adzuna and Jooble; and (7)~a zero-cost FAQ chatbot using Qdrant vector search with no LLM inference at query time. Implemented with Node.js~18, React~18, and PostgreSQL~14, the platform delivers a production-ready recruitment ecosystem bridging the gap between candidate preparation and employer evaluation.
\end{abstract}

\begin{IEEEkeywords}
multi-agent systems, recruitment automation, TF-IDF, cosine similarity, explainable AI, interview scheduling, coding assessment, career roadmap, large language models, job recommendation
\end{IEEEkeywords}

% ─────────────────────────────────────────────────────────────
\section{Introduction}
% ─────────────────────────────────────────────────────────────

Modern recruitment is characterised by exponentially growing application volumes and increasingly complex skill requirements. Studies indicate that recruiters spend an average of 23 hours screening resumes for a single hire~\cite{b1}, while approximately 75\% of qualified candidates are eliminated by inadequate automated screening systems~\cite{b2}. Most commercial Applicant Tracking Systems (ATS) rely on keyword matching that fails to capture semantic nuance and provide no transparent explanation to either recruiter or candidate~\cite{b3}.

Beyond screening, the hiring lifecycle suffers from siloed tooling: interview scheduling is manual and error-prone, coding skills are evaluated on separate third-party platforms, and resume improvement and career-planning tools are entirely disconnected from the application process. This fragmentation creates friction for both job seekers, who must navigate multiple platforms, and for recruiters, who must reconcile outputs from disconnected systems.

This paper introduces the \textbf{AI-Powered Agentic Hiring Platform}, a unified recruitment ecosystem that addresses all these gaps through a modular multi-agent architecture. Our key contributions are:

\begin{itemize}
\item A \textbf{hybrid shortlisting agent} that fuses TF-IDF vector similarity with Gemini~AI semantic reasoning, producing explainable skill-level match reports for every candidate.
\item A \textbf{break-aware round-robin scheduler} guaranteeing fair interviewer load distribution ($\Delta_{\text{load}}\!\leq\!1$) with $O(|C|)$ time complexity.
\item A \textbf{coding assessment module} supporting server-side execution in Python, C\texttt{++}, Java, and JavaScript~(Node.js) via the Piston API, enabling technical skill validation entirely within the platform.
\item An \textbf{AI resume optimiser} powered by Gemini with deterministic scoring, a hallucination guard, and downloadable PDF output.
\item An \textbf{AI career roadmap generator} driven by Groq Llama~3.3-70B and rendered as an interactive React-Flow metro-map.
\item A \textbf{profile-driven job recommender} performing per-skill parallel queries against Adzuna and Jooble APIs with automatic deduplication.
\end{itemize}

The remainder of this paper is organised as follows. Section~II reviews related work. Section~III presents the system architecture. Section~IV details core algorithms and methodologies. Section~V discusses experimental characteristics. Section~VI addresses limitations, and Section~VII concludes.

% ─────────────────────────────────────────────────────────────
\section{Related Work}
% ─────────────────────────────────────────────────────────────

\subsection{Recruitment Automation Systems}
Several commercial and academic systems have attempted to automate aspects of the recruitment process. LinkedIn Recruiter~\cite{b4} employs collaborative filtering to recommend candidates based on recruiter behaviour patterns; however, it lacks transparent explanation mechanisms and relies heavily on user network effects rather than objective skill matching. Research by Chen et al.~\cite{b5} proposed a deep learning-based resume screening system using convolutional neural networks, yet their approach treated resume analysis as a black-box classification problem without providing interpretable results to recruiters.

\subsection{Multi-Agent Systems in HR}
Multi-agent systems (MAS) have been explored for distributed problem-solving in various domains~\cite{b6}. In the recruitment context, Faliagka et al.~\cite{b7} developed a multi-criteria decision analysis framework for candidate ranking. However, their system required extensive manual configuration of criteria weights and lacked adaptive learning. Our work differs by implementing autonomous agents with specialised functions that collaborate through well-defined REST interfaces, enabling modular development and independent scaling without manual tuning.

\subsection{Interview Scheduling Optimisation}
Interview scheduling has been studied as a constraint satisfaction problem~\cite{b8}. Traditional approaches use genetic algorithms or simulated annealing to find optimal schedules, incurring high computational cost and offering no fairness guarantees. Our break-aware round-robin algorithm provides a provable near-perfect distribution guarantee while maintaining $O(|C|)$ computational complexity, a practical advantage for web-scale deployment.

\subsection{AI-Assisted Candidate Development}
LLM-based resume enhancement~\cite{b9,b10} and career path generation have been explored in isolation. Natural language processing for structured resume parsing has been examined in~\cite{b11}, yet no prior work integrates ATS optimisation, interactive career roadmapping, and automated coding assessment within a single unified recruitment platform.

% ─────────────────────────────────────────────────────────────
\section{System Architecture}
% ─────────────────────────────────────────────────────────────

\subsection{Three-Tier Design}

The platform follows a three-tier architecture comprising a Presentation Layer, an Application Layer, and a Data Layer, as illustrated in Fig.~\ref{fig:architecture}.

\begin{figure}[htbp]
\centerline{\includegraphics[width=\columnwidth]{system_architecture.png}}
\caption{Multi-agent system architecture showing three-tier design with six specialised AI agents, external API integrations, and the PostgreSQL data layer.}
\label{fig:architecture}
\end{figure}

The \textbf{Presentation Layer} consists of React~18 (Vite build tool) frontend modules with Tailwind CSS styling and React-Router~v6 navigation, serving distinct UIs for job seekers, recruiters, and administrators. The \textbf{Application Layer} houses six autonomous AI agents, RESTful API services built on Express.js, and JWT/bcrypt authentication middleware implementing role-based access control (RBAC). The \textbf{Data Layer} integrates PostgreSQL~14 (hosted on Neon cloud) with 30+ normalised tables and interfaces with four external services: Adzuna API, Jooble API, Groq API, and Google Gemini API.

\subsection{Multi-Agent System Design}

Six specialised autonomous agents each handle a distinct domain:

\begin{enumerate}
\item \textbf{Shortlist Agent:} Hybrid TF-IDF cosine similarity + Gemini~AI candidate ranking with explainability output.
\item \textbf{Scheduler Agent:} Break-aware round-robin interview time-slot assignment using Fisher-Yates shuffling.
\item \textbf{Aggregator Agent:} Parallel job consolidation from Adzuna, Jooble, and the internal database with deduplication.
\item \textbf{Optimizer Agent:} Gemini-powered ATS resume enhancement with deterministic pre/post scoring and hallucination detection.
\item \textbf{Roadmap Agent:} Groq Llama~3.3-70B career learning-path generation formatted as a React-Flow node graph.
\item \textbf{Assessor Agent:} Piston-backed remote code execution and automated test-case evaluation in four languages.
\item \textbf{Chatbot Agent:} Zero-cost semantic FAQ retrieval using Qdrant vector search — no LLM call at query time.
\end{enumerate}

\subsection{Database Schema}

The PostgreSQL schema is partitioned into functional domains:
\begin{itemize}
\item \textit{User \& Auth:} \texttt{users}, \texttt{credentials} (Google OAuth supported).
\item \textit{Candidate:} \texttt{candidates}, \texttt{candidate\_education}, \texttt{candidate\_resumes}.
\item \textit{Job Management:} \texttt{job\_postings}, \texttt{job\_requirements}.
\item \textit{Applications:} \texttt{job\_applications}, \texttt{job\_application\_answers}.
\item \textit{Assessments:} \texttt{tests}, \texttt{test\_questions}, \texttt{test\_attempts}.
\item \textit{AI Results:} \texttt{interviews}, \texttt{notifications}, \texttt{optimized\_resumes}.
\end{itemize}
Referential integrity is enforced via cascading foreign keys; soft deletion on job postings preserves application history.

% ─────────────────────────────────────────────────────────────
\section{Core Algorithms and Methodology}
% ─────────────────────────────────────────────────────────────

\subsection{Hybrid AI Candidate Shortlisting}

The Shortlist Agent implements a two-phase pipeline. In \textit{Phase~1}, TF-IDF vectorisation produces a deterministic, auditable numeric score. In \textit{Phase~2}, Gemini~1.5-Flash performs semantic reasoning to capture synonymy and contextual relevance missed by exact-term matching.

\subsubsection{Text Normalisation and Skill Aliasing}
Before vectorisation, a curated alias map resolves common abbreviations: \texttt{js}$\to$\texttt{javascript}, \texttt{node.js}$\to$\texttt{nodejs}, \texttt{k8s}$\to$\texttt{kubernetes}, \texttt{ml}$\to$\texttt{machine learning}, \texttt{postgres}$\to$\texttt{postgresql}, etc. Aliases are sorted by descending length to prevent partial matches. Text undergoes tokenisation, lowercasing, punctuation removal, and stopword filtering using the \texttt{natural} NLP library.

\subsubsection{TF-IDF Contextual Vectorisation}
An enriched candidate context $D_c$ is constructed for each candidate:
\begin{equation}
D_c = S_{\text{skills}} \oplus E_{\text{education}} \oplus R_{\text{resume}}
\end{equation}
where $\oplus$ denotes concatenation. Structuring data this way front-loads important keywords, boosting their term frequency in the two-document corpus $\{D_j, D_c\}$ built from the job description $D_j$. The TF-IDF weight for term $t$ in document $d$ is:
\begin{equation}
w_{t,d} = \mathrm{tf}(t,d) \times \log\!\left(\frac{N}{\mathrm{df}(t)}\right), \quad N{=}2
\end{equation}

\subsubsection{Cosine Similarity Scoring}
The final match score is derived from the cosine angle between job vector $\vec{V}_j$ and candidate vector $\vec{V}_c$:
\begin{equation}
S = \frac{\vec{V}_j \cdot \vec{V}_c}{\|\vec{V}_j\|\,\|\vec{V}_c\|} \times 100 \quad \in [0, 100]
\end{equation}

\subsubsection{Explainability and Degree Hierarchy}
A degree hierarchy maps qualifications to ordinal levels (Certificate$=1$, Bachelor$=2$, Masters$=3$, PhD$=4$) to evaluate education compatibility. The agent returns per-candidate: matched skills list, missing skills list, education compatibility flag, seniority tier alignment, and a natural-language summary string (e.g., ``8 of 10 required skills matched. Missing: Docker, Kubernetes.'').

\begin{figure}[htbp]
\centerline{\includegraphics[width=\columnwidth]{shortlisting_flowchart.png}}
\caption{AI-driven candidate shortlisting pipeline. Resumes undergo text normalisation and TF-IDF vectorisation followed by cosine similarity scoring. Gemini~AI adds semantic reasoning for a final explainable score with matched/missing skill breakdown.}
\label{fig:shortlist}
\end{figure}

\subsection{Break-Aware Round-Robin Scheduling}

The Scheduler Agent distributes $|C|$ candidates across $|I|$ interviewers in $O(|I|+|C|)$ time.

\begin{algorithm}[htbp]
\caption{Break-Aware Round-Robin Scheduling}
\label{alg:schedule}
\begin{algorithmic}[1]
\REQUIRE Candidates $C$, Interviewers $I$, start $t_0$, slot duration $d$, break freq $k$, break dur $B$
\ENSURE Schedule $\mathcal{S}$
\STATE $I \leftarrow \text{FisherYatesShuffle}(I)$
\STATE $T[i]\!\leftarrow\!t_0$,\; $cnt[i]\!\leftarrow\!0$ \; $\forall\, i\!\in\! I$
\FOR{$j = 0$ \TO $|C|-1$}
  \STATE $idx \leftarrow j \bmod |I|$
  \STATE $\mathcal{S}.\mathrm{add}\bigl(C[j],\;I[idx],\;T[idx],\;T[idx]+d\bigr)$
  \STATE $T[idx] \mathrel{+}= d$;\quad $cnt[idx] \mathrel{+}= 1$
  \IF{$cnt[idx] \bmod k = 0$ \AND $j < |C|-1$}
    \STATE $T[idx] \mathrel{+}= B$ \COMMENT{insert mandatory break}
  \ENDIF
\ENDFOR
\RETURN $\mathcal{S}$
\end{algorithmic}
\end{algorithm}

\textbf{Load Balance Guarantee.} The maximum interview-count difference between any two interviewers is bounded by:
\begin{equation}
\Delta_{\text{load}} = \left\lceil \frac{|C|}{|I|} \right\rceil - \left\lfloor \frac{|C|}{|I|} \right\rfloor \leq 1
\end{equation}

Fisher-Yates shuffling eliminates alphabetical bias in initial assignment order, ensuring that different recruiters are not systematically favoured across multiple scheduling runs.

\begin{figure}[htbp]
\centerline{\includegraphics[width=\columnwidth]{scheduling_algorithm.png}}
\caption{Break-aware round-robin scheduling flowchart. Fisher-Yates shuffle randomises interviewer order; modulo arithmetic distributes candidates; automatic break insertion prevents evaluator fatigue. The algorithm guarantees $\Delta_{\text{load}} \leq 1$.}
\label{fig:scheduling}
\end{figure}

\subsection{Coding Assessment Module}

The Assessor Agent provides an end-to-end technical evaluation environment:

\textbf{Recruiter Interface.} Recruiters create tests with multiple-choice and programming questions, attach hidden test cases with expected outputs, configure per-question time limits, and assign tests to specific job postings.

\textbf{Candidate Interface.} Candidates access an in-browser Monaco editor with syntax highlighting, receive real-time compilation feedback, and submit solutions before time expiry.

\textbf{Execution Backend.} Source code is dispatched to the hosted Piston API (\texttt{emkc.org/api/v2/piston}) supporting Python\,3, C\texttt{++}, Java, and JavaScript~(Node.js). A concurrency limiter (max 3 parallel requests) and a 300\,ms inter-request delay respect Piston rate limits. Output normalisation (CRLF$\to$LF, trailing whitespace stripping, internal whitespace collapsing) ensures robust expected-vs-actual comparison:
\begin{equation}
\text{Score} = \left\lfloor \frac{P}{T} \times 100 \right\rfloor \%
\end{equation}
where $P$ = test cases passed and $T$ = total test cases.

\subsection{AI Resume Optimiser}

The Optimizer Agent processes PDF or DOCX uploads using \texttt{pdf-parse} and \texttt{mammoth} respectively. The optimisation pipeline is:
\begin{enumerate}
\item \textit{Deterministic baseline scoring}: keyword density, action-verb detection, and quantification analysis produce a pre-optimisation ATS score.
\item \textit{Gemini enhancement}: a structured prompt (temperature 0.4) instructs Gemini~Flash to rewrite bullet points with strong action verbs, integrate missing keywords naturally, and improve quantification—without inventing companies, dates, or degrees.
\item \textit{Hallucination guard}: the optimised output is discarded if new year tokens (not present in the original) appear, or if output length drops below 50\% of input length.
\item \textit{Re-scoring and export}: post-optimisation ATS score is computed; delta is provided to the candidate; a downloadable PDF is generated and persisted to \texttt{optimized\_resumes}.
\end{enumerate}

\subsection{AI Career Roadmap Generator}

The Roadmap Agent accepts a target skill and current proficiency level (Beginner / Intermediate / Advanced). It queries Groq Llama~3.3-70B-Versatile (temperature 0.3, JSON mode) for a structured learning graph with 6--10 milestone nodes, each containing curated resource links (Coursera, Udemy, official documentation) and estimated time-to-completion. The JSON output comprising \texttt{nodes} (id, label, resources, estimated\_time, position) and \texttt{edges} (source, target) is consumed directly by React-Flow to render an interactive metro-map diagram in the candidate dashboard.

\subsection{Profile-Driven Job Recommender}

The Recommender reads up to five skills and the candidate's job title from the \texttt{candidates} table, then issues \emph{per-skill} parallel API queries to Jooble and Adzuna. Results are merged and deduplicated by a composite key of normalised title + company, returning the top~30 unique listings with source attribution and matched-skill tags.

\subsection{Zero-Cost FAQ Chatbot via Qdrant}

The Chatbot Agent answers platform-related user queries without invoking any LLM at query time, eliminating per-token API costs and reducing latency to under 100\,ms.

\textbf{Offline Indexing (one-time).} A corpus of 50 curated FAQ entries (question, answer, keywords) is encoded using the \texttt{all-MiniLM-L6-v2} Sentence Transformer model (384-dimensional vectors). These embeddings are bulk-uploaded to a Qdrant cloud collection (\texttt{faq\_collection}, hosted on GCP europe-west3) using cosine distance metric:
\begin{equation}
\text{sim}(q, d) = \frac{\vec{q}\cdot\vec{d}}{\|\vec{q}\|\,\|\vec{d}\|}
\end{equation}
All payloads (question, answer, category, keywords) are stored inline within each Qdrant point.

\textbf{Online Retrieval (per query).} When a user submits a natural language query, the Node.js backend encodes it with the same MiniLM model and calls \texttt{qdrant\_client.query\_points()} with \texttt{top\_k=3}. The three highest-cosine-similarity FAQ matches are returned instantly. No generative model is called; the answer is the stored FAQ 	exttt{answer} field of the top-ranked result. This design guarantees deterministic, hallucination-free responses grounded entirely in curated platform documentation.

The FAQ corpus spans 9 categories: Account Management, Job Applications, Job Posting, Search \& Filters, Profile Features, Notifications, Payments, Platform Features, and Troubleshooting.

% ─────────────────────────────────────────────────────────────
\section{System Implementation and Results}
% ─────────────────────────────────────────────────────────────

The platform is implemented using Node.js~18 and Express.js for the backend, React~18 (Vite + Tailwind CSS) for the frontend, and PostgreSQL~14 (Neon cloud) for data persistence. Authentication is handled via JWT + bcrypt with optional Google OAuth. AI capabilities are provided by Groq Llama~3.3-70B-Versatile (career roadmap), Google Gemini~1.5-Flash (shortlisting, resume optimisation, and resume parsing), and the Piston API for remote code execution. The \texttt{natural} Node.js NLP library handles TF-IDF vectorisation and tokenisation. Fig.~\ref{fig:dataflow} shows the complete data flow through the system.

\begin{figure}[htbp]
\centerline{\includegraphics[width=\columnwidth]{data_flow_architecture.png}}
\caption{Data flow architecture. Job seekers and recruiters interact through the React frontend; stateless AI agents process data and persist results to PostgreSQL; Adzuna, Jooble, Groq, and Gemini are integrated at the application layer.}
\label{fig:dataflow}
\end{figure}

\subsection{AI Shortlisting Algorithm Properties}

\begin{table}[htbp]
\caption{AI Shortlisting Algorithm Characteristics}
\label{tab:shortlist_props}
\begin{center}
\begin{tabular}{|l|l|}
\hline
\textbf{Property} & \textbf{Value} \\
\hline
Vectorisation Library    & \texttt{natural} (Node.js NLP) \\
AI Semantic Layer        & Google Gemini 1.5-Flash \\
Score Range              & 0--100 integer \% \\
Approx.\ Processing Time & $\sim$45\,ms per resume (TF-IDF phase) \\
TF-IDF Corpus Size       & Dynamic (2 documents per candidate) \\
Time Complexity          & $O(c \times n^2)$, $c$=candidates, $n$=terms \\
Explainability Output    & Matched skills, missing skills, \\
                         & education fit, seniority alignment \\
\hline
\end{tabular}
\end{center}
\end{table}

\subsection{Scheduling Algorithm Properties}

\begin{table}[htbp]
\caption{Scheduling Algorithm Properties}
\label{tab:sched_props}
\begin{center}
\begin{tabular}{|l|l|}
\hline
\textbf{Property} & \textbf{Value} \\
\hline
Initial Order         & Fisher-Yates shuffle ($O(|I|)$) \\
Assignment Strategy   & Round-robin: $idx = j \bmod |I|$ \\
Load Guarantee        & $\Delta_{\text{load}} \leq 1$ \\
Time Complexity       & $O(|I| + |C|) \approx O(|C|)$ \\
Break Insertion       & Per-interviewer, configurable $k$ \\
Determinism           & Guaranteed post-shuffle \\
State Tracking        & Per-interviewer time + count \\
\hline
\end{tabular}
\end{center}
\end{table}

\subsection{Core API Endpoints}

Table~\ref{tab:api} maps each platform capability to its REST endpoint and responsible agent.

\begin{table}[htbp]
\caption{Core API Endpoints and Responsible Agents}
\label{tab:api}
\begin{center}
\resizebox{\columnwidth}{!}{%
\begin{tabular}{|l|c|l|}
\hline
\textbf{Endpoint Path} & \textbf{Method} & \textbf{Agent} \\
\hline
\texttt{/ai-tools/shortlist/:jobId}       & POST & Shortlist \\
\texttt{/interviews/auto-schedule/:jobId} & POST & Scheduler \\
\texttt{/jobs/aggregated}                 & GET  & Aggregator \\
\texttt{/ai-resume/optimize}              & POST & Optimizer \\
\texttt{/career-roadmap/generate}         & POST & Roadmap \\
\texttt{/coding/submit}                   & POST & Assessor \\
\texttt{/recommended-jobs}                & GET  & Recommender \\
\texttt{/chatbot/query}                   & POST & Chatbot (Qdrant) \\
\hline
\end{tabular}%
}
\end{center}
\end{table}

% ─────────────────────────────────────────────────────────────
\section{Discussion}
% ─────────────────────────────────────────────────────────────

\subsection{Explainability, Fairness, and Trust}

The TF-IDF cosine similarity layer provides full audit traceability—every term weight is inspectable—satisfying GDPR ``right to explanation'' requirements. The Gemini semantic layer adds depth without sacrificing transparency since TF-IDF scores remain the primary ranking signal. The Fisher-Yates shuffle prevents alphabetical bias in interviewer assignment, and the $\Delta_{\text{load}} \leq 1$ guarantee ensures equitable workload distribution. Configurable break parameters ($k$, $B$) allow administrators to enforce evaluator quality standards based on organisational norms. By consolidating resume optimisation, roadmapping, coding tests, and job recommendations into one authenticated ecosystem, candidates gain a continuous improvement loop and recruiters gain richer signal about candidate initiative.

\subsection{Limitations and Future Work}

\begin{itemize}
\item TF-IDF does not capture deep synonymy; augmenting with Sentence-BERT embeddings would improve semantic recall.
\item Piston API rate limits (1~req/200\,ms) constrain throughput; a self-hosted instance would remove this bottleneck.
\item The hallucination guard uses heuristic rules; an NLI-based faithfulness check would be more robust.
\item Future work includes transformer-based resume matching, reinforcement learning for adaptive weight tuning, video sentiment analysis, and multilingual resume support.
\end{itemize}

% ─────────────────────────────────────────────────────────────
\section{Conclusion}
% ─────────────────────────────────────────────────────────────

This paper presented the AI-Powered Agentic Hiring Platform, a unified recruitment ecosystem addressing fragmentation across the hiring lifecycle through six autonomous agents. The hybrid TF-IDF + Gemini shortlisting pipeline delivers explainable candidate evaluation in $\sim$45\,ms per resume. The break-aware round-robin algorithm guarantees $\Delta_{\text{load}} \leq 1$ load balance in $O(|C|)$ time. Piston-backed coding assessment, Gemini resume optimisation, Groq Llama~3.3 career roadmapping, and profile-driven job recommendation complete an integrated candidate development loop—all within a single production-ready platform built on Node.js~18, React~18, and PostgreSQL~14.

\section*{Acknowledgment}

The authors thank the Department of Artificial Intelligence and Data Science, Viswajyothi College of Engineering and Technology, for computational resources and research support.

% ─────────────────────────────────────────────────────────────
\begin{thebibliography}{00}

\bibitem{b1} M. Johnson and R. Smith, ``Time analysis of modern recruitment processes: A quantitative study,'' \textit{Journal of Human Resource Management}, vol.~34, no.~2, pp.~145--163, 2024.

\bibitem{b2} A. Williams, ``The candidate filtering paradox: How ATS systems eliminate qualified applicants,'' \textit{HR Technology Quarterly}, vol.~12, no.~4, pp.~78--92, 2023.

\bibitem{b3} K. Zhang, L. Chen, and M. Wang, ``Machine learning applications in human resource management: A systematic review,'' \textit{IEEE Access}, vol.~11, pp.~45678--45702, 2023.

\bibitem{b4} LinkedIn Corporation, ``LinkedIn Recruiter: Technical white paper on recommendation algorithms,'' Technical Report, LinkedIn Engineering, 2023.

\bibitem{b5} Y. Chen, X. Liu, and J. Wang, ``Deep learning-based resume screening using convolutional neural networks,'' in \textit{Proc.\ IEEE Int.\ Conf.\ on Data Mining (ICDM)}, Beijing, China, 2023, pp.~234--243.

\bibitem{b6} M. Wooldridge, \textit{An Introduction to MultiAgent Systems}, 2nd~ed. Chichester, UK: Wiley, 2009.

\bibitem{b7} E. Faliagka, A. Tsakalidis, and G. Tzimas, ``An integrated e-recruitment system for automated personality mining and applicant ranking,'' \textit{Internet Research}, vol.~22, no.~5, pp.~551--568, 2012.

\bibitem{b8} R. Kumar and S. Patel, ``Optimization techniques for interview scheduling: A comparative analysis,'' in \textit{Proc.\ Int.\ Conf.\ on Artificial Intelligence and Applications (ICAIA)}, Chennai, India, 2022, pp.~567--574.

\bibitem{b9} J. Devlin, M. Chang, K. Lee, and K. Toutanova, ``BERT: Pre-training of deep bidirectional transformers for language understanding,'' in \textit{Proc.\ NAACL-HLT}, Minneapolis, MN, USA, 2019, pp.~4171--4186.

\bibitem{b10} T. Brown et al., ``Language models are few-shot learners,'' in \textit{Advances in Neural Information Processing Systems (NeurIPS)}, vol.~33, 2020, pp.~1877--1901.

\end{thebibliography}

\end{document}
